\begin{Solbook}{1}
 \par \noindent  a) $7\sqrt {3}+9$, \hspace {2cm} b) $\dfrac {\sqrt {x}-1}{2}$  \hfill \qedEX 
\end{Solbook}
\begin{Solbook}{2}
 \par \noindent  a) $x=2$, \hspace {2cm} b) $x\ge \dfrac {5}{2}$  \hfill \qedEX 
\end{Solbook}
\begin{Solbook}{3}
 \par \noindent  Ban tổ chức đã bán được 250 vé loại I và 325 vé loại II.  \hfill \qedEX 
\end{Solbook}
\begin{Solbook}{4}
 \par \noindent  $AC=250.\tan {38^{\circ }}$. Từ đó $EC=250.\tan {38^{\circ }}+1,5\approx 197 $ (m).  \hfill \qedEX 
\end{Solbook}
\begin{Solbook}{5}
 \par \noindent  \immini { \begin {enumerate} \item Do $BC\perp AC$ nên $OM\perp AC$. \par Chứng minh $\triangle OAM=\triangle OCM$ suy ra $\widehat {OCM}=90^\circ $. \par \item $\triangle MAH \backsim \triangle ABC$ (g.g) suy ra $\dfrac {MA}{AB}=\dfrac {MH}{AC}$. \par $\triangle MAI \backsim \triangle ABH$ (g.g) suy ra $\dfrac {MA}{AB}=\dfrac {MI}{HA}$. \par Suy ra $\dfrac {MI}{HA}=\dfrac {MH}{AC}$ nên $\dfrac {MI}{MH}=\dfrac {HA}{AC}=\dfrac {1}{2}$. \par Vậy $I$ là trung điểm $MH$, suy ra $MI=IH$. \end {enumerate} }{ \begin {tikzpicture}[scale=1, font=\footnotesize , >=stealth] \def \R {2} \coordinate (O) at (0,0); \coordinate (A) at (-\R ,0); \coordinate (B) at (\R ,0); \coordinate (C) at (60:\R ); \coordinate (H) at ($(A)!(O)!(C)$); \coordinate (F) at ($(H)!(A)!(B)$); \coordinate (x) at ($(A)!1cm!90:(O)$); \coordinate (M) at (intersection of A--x and O--H); \draw (-2,0) arc[start angle=180, end angle=0, radius=2]; \draw (A)--(B)--(C)--cycle; \draw (O)--(M)--(C)--cycle (M)--(A) ; \draw (O)--(H) (A)--(I); \coordinate (I) at ($(M)!0.5!(H)$); \draw [dashed] (B)--(F); \foreach \p /\pos in {A/left,B/right,O/below,C/above right,M/above,H/above,I/left,F/above left} {\filldraw [fill=white] (\p ) circle(1.5pt) node[\pos ] {$\p $}; } \end {tikzpicture} } \hfill \qedEX 
\end{Solbook}
\begin{Solbook}{6}
 \par \noindent  $x^{2}-2x-12=\sqrt {2x-1}$ \par $x^2-2x-15=\sqrt {2x-1}-3$ \par $(x-5)(x+3)=\dfrac {2(x-5)}{\sqrt {2x-1}+3}$ \par Đưa về phương trình tích. Phương trình có nghiệm duy nhất $x=5$.  \hfill \qedEX 
\end{Solbook}
